\chapter{Analyse des besoins}
\label{chap:analyse}

\section{Démarche d'analyse}
L'analyse des besoins a été conduite à partir des usages réels du parc informatique, puis traduite en services, rôles, règles de gestion et critères d'acceptation. Cette démarche a permis de relier chaque fonctionnalité à un besoin métier observable: inventorier les actifs, contrôler les accès, tracer les mouvements, suivre la maintenance et disposer d'une vision de pilotage fiable.

La logique adoptée repose sur trois principes:
\begin{itemize}
  \item partir des processus métier avant les écrans;
  \item relier chaque besoin à un acteur et à une responsabilité explicite;
  \item distinguer clairement fonctions, contraintes et règles de gouvernance.
\end{itemize}

\section{Analyse fonctionnelle détaillée}
La plateforme peut être décomposée en cinq chaînes de valeur:
\begin{enumerate}
  \item \textbf{chaîne d'accès}: inscription, vérification d'email, approbation administrateur, authentification et contrôle d'accès;
  \item \textbf{chaîne d'inventaire}: création, mise à jour, classification, illustration photo et archivage documentaire des équipements;
  \item \textbf{chaîne d'affectation}: demande de matériel par l'employé, réservation, validation, activation, restitution et historisation d'un matériel;
  \item \textbf{chaîne de maintenance}: déclaration d'incident, création d'intervention, suivi kanban, clôture et remise en disponibilité;
  \item \textbf{chaîne de pilotage}: dashboard, analytics, notifications, rapports et audit.
\end{enumerate}

Cette structuration met en évidence que le projet ne se limite pas à une liste d'équipements. Il couvre un cycle de vie complet, depuis l'entrée d'un actif dans l'inventaire jusqu'à son affectation, sa maintenance, son retour et sa traçabilité.

\section{Identification des acteurs}
Les acteurs principaux du système sont présentes dans le tableau \ref{tab:acteurs}.

\begin{table}[H]
\centering
\caption{Acteurs du système}
\label{tab:acteurs}
\begin{tabularx}{\textwidth}{>{\bfseries}p{3.3cm}X}
\toprule
Acteur & Rôle dans le système \\
\midrule
Administrateur & Gère les comptes, approuve les inscriptions, administre les rôles et permissions, consulte les audits et supervise l'ensemble du système. \\
Responsable informatique & Gère l'inventaire, pilote les affectations, suit les indicateurs et coordonne l'activité opérationnelle du parc. \\
Technicien & Prend en charge les maintenances, met à jour les interventions, suit les tickets et fait progresser le kanban technique. \\
Employé / bénéficiaire & Consulte ses affectations, demande un nouveau matériel, reçoit les notifications et peut déclarer un incident sur ses propres équipements. \\
Système & Exécute les traitements transverses: JWT, audit, notifications, persistance, WebSocket et génération documentaire. \\
\bottomrule
\end{tabularx}
\end{table}

Dans l'implémentation finale, le terme \enquote{bénéficiaire} demeure dans certains objets d'affectation pour désigner le destinataire d'un matériel. En revanche, le rôle officiel exposé dans la gestion des comptes et l'inscription est \enquote{Employé / bénéficiaire}. Ce choix évite la confusion entre vocabulaire métier et vocabulaire de gestion des utilisateurs.

\section{Description des cas d'utilisation}
Les principaux cas d'utilisation identifiés sont:
\begin{enumerate}
  \item s'inscrire avec un rôle métier autorisé;
  \item vérifier son adresse e-mail;
  \item approuver ou rejeter un compte depuis l'administration;
  \item se connecter à la plateforme;
  \item consulter le tableau de bord;
  \item gérer les utilisateurs et leurs droits;
  \item créer et mettre à jour des équipements;
  \item téléverser une photo réelle et un document PDF pour un équipement;
  \item déposer une demande de matériel depuis l'espace employé bénéficiaire;
  \item traiter une demande de matériel par validation ou rejet côté administration;
  \item créer, valider et restituer une affectation;
  \item déclarer et suivre une maintenance;
  \item consulter les notifications et les rapports.
\end{enumerate}

\section{User stories}
Le backlog fonctionnel peut être exprimé sous la forme des user stories du tableau \ref{tab:user-stories}.

\begin{longtable}{p{1.5cm}p{4.8cm}p{3cm}p{5.1cm}}
\caption{Sélection de user stories}\\
\toprule
\textbf{ID} & \textbf{User story} & \textbf{Priorité} & \textbf{Critère d'acceptation principal} \\
\midrule
\endfirsthead
\toprule
\textbf{ID} & \textbf{User story} & \textbf{Priorité} & \textbf{Critère d'acceptation principal} \\
\midrule
\endhead
US01 & En tant qu'utilisateur, je veux m'inscrire avec mon rôle métier afin d'entrer dans le circuit de validation. & Haute & Le compte est créé en attente de vérification puis d'approbation. \\
US02 & En tant qu'administrateur, je veux approuver les comptes afin de contrôler qui accede au système. & Haute & Le compte devient actif seulement après vérification d'email et approbation. \\
US03 & En tant que responsable informatique, je veux créer un équipement afin d'enrichir l'inventaire centralisé. & Haute & L'équipement est persisté avec un code inventaire unique et un statut cohérent. \\
US04 & En tant que responsable informatique, je veux téléverser une photo et un PDF afin de documenter réellement l'actif. & Moyenne & Les fichiers sont stockés puis liés à la fiche équipement. \\
US05 & En tant qu'employé bénéficiaire, je veux consulter mes affectations afin de savoir quel matériel m'est rattaché. & Haute & Seules les affectations du compte connecté sont visibles. \\
US06 & En tant qu'administrateur ou responsable, je veux affecter un matériel à un employé éligible afin de fiabiliser le mouvement. & Haute & L'affectation est refusée si le bénéficiaire n'est pas actif, vérifié et approuvé. \\
US07 & En tant qu'employé, je veux déclarer une panne sur mon équipement afin de lancer une prise en charge. & Haute & La demande est refusée si l'équipement ne lui appartient pas ou n'est pas activement affecté. \\
US08 & En tant que technicien, je veux suivre les maintenances afin de prioriser les interventions. & Haute & Le tableau de maintenance et le kanban affichent les statuts et les dates clés. \\
US09 & En tant qu'administrateur, je veux observer les comptes, leurs statuts et leur activité afin de gouverner les accès. & Haute & Le dashboard utilisateurs présente recherche, filtres, statistiques et actions. \\
US10 & En tant qu'employé bénéficiaire, je veux demander un matériel afin d'exprimer un besoin avant affectation. & Haute & La demande est enregistrée en \texttt{PENDING} et visible par l'administrateur ou le responsable informatique. \\
US11 & En tant qu'administrateur ou responsable, je veux valider ou rejeter une demande de matériel afin de piloter les dotations. & Haute & La décision met à jour le statut et notifie le demandeur. \\
US12 & En tant que pilote du parc, je veux disposer d'un smoke test multi-rôles afin de valider rapidement les parcours critiques. & Moyenne & Le script de vérification confirme inscription, équipement, affectation, demande de matériel et maintenance. \\
\bottomrule
\label{tab:user-stories}
\end{longtable}

\section{Règles de gestion}
Les règles de gestion qui structurent le système sont les suivantes:
\begin{enumerate}
  \item un seul compte administrateur principal est créé automatiquement par le système;
  \item l'inscription frontend n'autorise jamais la création d'un administrateur;
  \item tout nouvel inscrit métier passe par les états \texttt{PENDING\_VERIFICATION} puis \texttt{PENDING\_APPROVAL};
  \item un équipement ne peut pas posséder plusieurs affectations actives simultanées;
  \item un équipement en maintenance n'est pas disponible pour une nouvelle affectation;
  \item une affectation ne peut viser qu'un employé bénéficiaire actif, vérifié et approuvé;
  \item une demande de matériel créée par un employé démarre toujours à l'état \texttt{PENDING};
  \item seules les personnes habilitées à gérer les affectations peuvent approuver, rejeter ou clôturer une demande de matériel;
  \item un employé ne peut déclarer une maintenance que sur un équipement qui lui est activement affecté;
  \item toute action sensible doit être historisée dans les journaux d'activité et d'audit;
  \item les listes de référence ne doivent proposer que des utilisateurs éligibles aux workflows cibles.
\end{enumerate}

\section{Spécifications fonctionnelles}
Les spécifications fonctionnelles sont résumées dans le tableau \ref{tab:specs-fonctionnelles}.

\begin{table}[H]
\centering
\caption{Spécifications fonctionnelles synthétiques}
\label{tab:specs-fonctionnelles}
\begin{tabularx}{\textwidth}{>{\bfseries}p{3cm}X}
\toprule
Module & Fonctions couvertes \\
\midrule
Authentification & connexion, inscription avec choix de rôle, vérification d'e-mail, logout, refresh token, oubli de mot de passe \\
Administration utilisateurs & approbation, rejet, suspension, activation, désactivation, changement de rôle, recherche et supervision \\
Dashboard & indicateurs globaux, activité récente, synthèse du parc, maintenance et notifications \\
Équipements & consultation, création, mise à jour, upload photo, upload document PDF, lecture des statuts et garanties \\
Affectations & demande de matériel employé, traitement des demandes, création, détail, approbation, restitution, historique et contrôle des éligibilités \\
Maintenance & création, suivi, changement de statut, vue kanban, vérification des droits et des propriétés métier \\
Notifications & consultation du flux, marquage en lecture, diffusion temps réel \\
Rapports & export PDF/XLSX et restitution des données de pilotage \\
\bottomrule
\end{tabularx}
\end{table}

\section{Spécifications non fonctionnelles}
Les exigences non fonctionnelles suivantes ont structuré la solution:
\begin{table}[H]
\centering
\caption{Spécifications non fonctionnelles}
\label{tab:specs-non-fonctionnelles}
\begin{tabularx}{\textwidth}{>{\bfseries}p{3.3cm}X}
\toprule
Catégorie & Exigence \\
\midrule
Sécurité & JWT, hachage BCrypt, routes protégées, RBAC, en-têtes de sécurité, journalisation \\
Fiabilité & transactions cohérentes sur les affectations et maintenances, vérifications métier côté serveur \\
Maintenabilité & séparation frontend/backend, services clairs, DTOs, composants React réutilisables \\
Portabilité & exécution locale, build React intégré à Spring Boot, Docker Compose disponible \\
Ergonomie & interface SaaS claire, responsive, lisible sur mobile, tablette et desktop \\
Traçabilité & audit logs, activity logs, notifications et historique des changements \\
Évolutivité & architecture modulaire permettant l'ajout de nouveaux modules et de nouveaux tests \\
\bottomrule
\end{tabularx}
\end{table}

\section{Synthèse du chapitre}
L'analyse des besoins montre que la plateforme traite une gouvernance complète du parc: accès, inventaire, demande de matériel, affectation, maintenance, notification et pilotage. Elle met également en évidence trois exigences majeures qui ont guidé les dernières corrections: la clarté du rôle \enquote{Employé / bénéficiaire}, la capacité pour l'utilisateur final d'exprimer un besoin de dotation, et la nécessité de faire respecter les workflows critiques par le backend lui-même, et non seulement par l'interface.









